\documentclass[11pt,a4paper]{article}
\usepackage[margin=25mm]{geometry}
\usepackage{amsmath,amssymb}
\usepackage{booktabs}
\usepackage{longtable}
\usepackage{array}
\usepackage[hidelinks]{hyperref}
\usepackage{xcolor}

\newcommand{\code}[1]{\texttt{#1}}
\newcommand{\psr}{B0329+54}
\setlength{\parskip}{0.5ex}

\title{Pulsar observing with the Acre Road SRT:\\
recording, folding, times of arrival and the data kept}
\author{Acre Road Observatory}
\date{30 September 2026}

\begin{document}
\maketitle

\begin{abstract}
The Acre Road 3\,m telescope detects PSR~\psr{} at 1.4\,GHz. It first did so on the night of 26--27 September 2026, at a matched signal-to-noise ratio (SNR) of 9.4 in 16\,h. PRESTO confirmed the detection independently at $9.9\sigma$.

This document describes the pulsar mode end to end:
\begin{itemize}
\item how the radio records a total-power series at 1\,ms;
\item how every sample is timed;
\item how the pulsar monitor schedules a run every day;
\item the barycentric phase model used to fold it;
\item how the fold is made and judged;
\item how a time of arrival (TOA) is measured and written for the timing program PINT;
\item which files hold what, and which of them are kept in git.
\end{itemize}
It closes with the checks that established each step and the questions still open. File and function names refer to \code{receiver\_scheduler/} in the repository.
\end{abstract}

\tableofcontents
\newpage

%-----------------------------------------------------------------------------
\section{Scope and the target}

The mode is built for PSR~\psr{} (J0332+5434) and nothing else. At 1.4\,GHz it is the only pulsar within reach of a 3\,m dish at the present system temperature: its mean flux density is $203\pm57$\,mJy, and the next candidate, B0950+08, would need of order twenty times the integration. Its parameters are therefore hard-wired in \code{pulsar\_fold.B0329}; there is no catalogue lookup and no pulsar picker.

\begin{table}[h]
\centering
\begin{tabular}{@{}ll@{}}
\toprule
Position (J2000, epoch MJD 56000) & 03:32:59.4096, +54:34:43.329 \\
Proper motion & $\mu_\alpha^* = 16.97$, $\mu_\delta = -10.37$ mas\,yr$^{-1}$ \\
$F_0$ at PEPOCH MJD 46473 (TDB) & 1.399541538720\,Hz ($P_0 = 0.714519699726$\,s) \\
$F_1$, $F_2$ & $-4.01197\times10^{-15}$\,s$^{-2}$, $5.3\times10^{-28}$\,s$^{-3}$ \\
DM & 26.7641\,pc\,cm$^{-3}$ \\
$S_{1400}$, $W_{50}$, $W_{10}$ & 203\,mJy, 6.6\,ms, 31--47\,ms \\
\bottomrule
\end{tabular}
\caption{The ephemeris used, from the ATNF catalogue (Hobbs et al.\ 2004).}
\end{table}

The pulse is not a single Gaussian. In the EPN profile at 1410\,MHz (von Hoensbroech \& Xilouris 1997) only 49\% of the pulse energy lies within $\pm5$\,ms of the main peak. A leading shoulder and two outriders, at about $-24$ and $+18$\,ms, carry the rest. The main peak is strongly circularly polarised, with $V/I = -0.34$. Both facts enter the expected sensitivity (\S\ref{sec:expect}).

%-----------------------------------------------------------------------------
\section{The instrument in pulsar mode}

\subsection{Radio and band}
The receiver is an Ettus USRP \textbf{B200}: one receive channel, AD9364 transceiver. It sits behind the feed, the SAWbird+ H1 front end (LNA and SAW filter) and the cable. Pulsar mode is not tied to the fixed H\,\textsc{i} instrument used for spectra, because it needs no spectral line. Its band is set by two configuration keys, which the scheduler passes to the receiver as environment variables:
\begin{itemize}
\item \code{pulsar\_sample\_rate\_hz}, as \code{H1\_PULSAR\_RATE}: default 32\,Msps;
\item \code{pulsar\_lo\_hz}, as \code{H1\_PULSAR\_LO}: default 1413\,MHz.
\end{itemize}
The band is therefore 1397--1429\,MHz.

A 64-channel test with the dish stowed found the whole receive chain flat to within 2\,dB over 1397.5--1428\,MHz, with every 0.5\,MHz channel at the radiometer noise level (1.00--1.04 times). Summing into a single channel then loses almost nothing:
\begin{itemize}
\item the effective bandwidth of the summed channel is 27.8\,MHz, against 28.2\,MHz for per-channel normalisation;
\item that compares with 7.4\,MHz effective for the fixed instrument's 8\,Msps;
\item SNR scales as $\sqrt{B}$, so the gain is a factor of 1.9.
\end{itemize}
The receiver therefore records the whole band as \textbf{one channel}. No passband correction is needed.

Two facts about the band edges:
\begin{itemize}
\item 1397--1400 and 1427--1429\,MHz lie outside the 1400--1427\,MHz radio-astronomy allocation. A single channel cannot drop them if interference appears there.
\item The AD9364's analogue filter is requested at twice the sample rate, capped at the chip's 56\,MHz. It only protects the ADC; anti-aliasing is done by the chip's digital decimation filters.
\end{itemize}

\subsection{Why one channel and no dedispersion}
The dispersion sweep across the band is
\begin{equation}
\Delta t = K\,\mathrm{DM}\,\left(f_\mathrm{lo}^{-2} - f_\mathrm{hi}^{-2}\right) \approx 2.5\,\mathrm{ms},
\qquad K = 4.148808\times10^{3}\ \mathrm{s\,MHz^2\,pc^{-1}\,cm^{3}}.
\end{equation}
Against the 7\,ms main peak ($\sigma \approx 3.0$\,ms) this smear costs about 1.5\% of the SNR. Coherent dedispersion, as was essential in the 408\,MHz system, would need an 80\,000-sample filter at 32\,Msps for that 1.5\%. Four channels would reduce the smear to 0.63\,ms and the loss to 0.1\%; this is noted for later, but not done.

\subsection{Signal path in the receiver}
\code{PulsarFlowgraph} in \code{b210\_h1\_receiver.py} is a GNU Radio flowgraph:
\[
\text{usrp\_source} \rightarrow |x|^2 \rightarrow \times\tfrac{1}{N} \rightarrow \text{sum of } N \rightarrow \text{row sink},
\qquad N = f_s\,\Delta t = 32\,000,
\]
giving rows of $\Delta t = 1$\,ms. Each row is the mean power of 32\,000 complex samples. With 16 channels (\code{H1\_PULSAR\_NCHAN}), a 16-point FFT filterbank replaces $|x|^2$. That option exists but is not the default: at 32\,Msps a 64-channel FFT overflowed the host.

\subsection{Clock, time and buffering}
\begin{description}
\item[Frequency.] The B200 runs from the Thunderbolt GPS receiver's 10\,MHz reference (\code{select\_clock\_source}). Lock is checked through the \code{ref\_locked} sensor rather than assumed. Within a run, sample timing is as good as the Thunderbolt.
\item[Absolute time.] \code{set\_device\_time} with \code{H1\_TIME\_SOURCE} = \code{auto} (the default):
  \begin{enumerate}
  \item switch the radio's time source to external and wait up to 2.5\,s for a PPS edge;
  \item just after an edge, set the next edge to be second $\lfloor t_\mathrm{host}\rfloor+1$;
  \item after 1.2\,s, check that the last PPS reads a whole second and that the radio agrees with the host to within the NTP error.
  \end{enumerate}
  With no PPS seen, the host's NTP clock is copied into the radio once (good to about 1--4\,ms), and the file says so. \code{external} refuses to record without a PPS; \code{internal} never looks for one. The whole second always comes from the host and only the sub-second from the PPS, so whether the Thunderbolt aligns its PPS to UTC or GPS seconds makes no difference: the two differ by a whole number of seconds. The B200's PPS input accepts 1.8--5\,V.
\item[USB buffering.] The radio is opened with \code{num\_recv\_frames=512} (\code{UHD\_DEVICE\_ARGS}). UHD's default of 16 frames holds only 1--2\,ms of samples at 32\,Msps, so any host stall longer than that dropped samples. Under the same 5-minute synthetic load, 16 frames gave 324 overflows and 512 frames gave none.
\end{description}

\subsection{Timing every row}
When the radio drops samples (an overflow), UHD tags the first sample after the gap with \code{rx\_time}, the device clock's reading for that sample. The stream's first sample is tagged the same way. These marks are what keep the time axis exact across gaps.

Downstream of the summing block, GNU Radio rounds a tag's offset to the nearest row: an offset of 499 samples lands in row 0 and 500 in row 1. The position of the gap within its row would be lost, and every row after it stamped up to $\pm0.5$\,ms wrong. The receiver therefore reads the tags \emph{before} the sum. \code{\_TagTap} uses GNU Radio's C++ \code{tag\_debug} block (silent, keeping every tag) on the radio's own output. A Python block doing the same work overflowed the radio every 9\,s at 32\,Msps. Each mark is stored as a fractional row, (sample index)/(samples per row), together with the device time.

The time of row $i$ is its \emph{middle}:
\begin{equation}
t_i = t_k + \Delta t\,\bigl(i + \tfrac12 - r_k\bigr),
\end{equation}
where $(r_k, t_k)$ is the latest mark whose row starts at or before row $i$. A stretch begins at the first row lying wholly after its mark; the row the gap falls in began on the previous mark. This is \code{pulsar\_fold.row\_times}, and \code{\_block\_times} applies the same rule block by block. Stamping rows at their start would make every TOA 0.5\,ms early.

%-----------------------------------------------------------------------------
\section{The recording}
\label{sec:file}

Each run writes one HDF5 file, \code{data/observations/YYYYMMDD\_HHMMSS\_pulsar.h5}. It is readable while being written (SWMR). Four hours at one channel are 58\,MB; sixteen hours are 230\,MB.

\begin{longtable}{@{}p{0.27\linewidth}p{0.68\linewidth}@{}}
\toprule
Dataset & Contents \\
\midrule
\code{power} & $N_\mathrm{rows}\times n_\mathrm{chan}$, float32: mean power per row in counts. \\
\code{frequency\_hz} & the channel centre(s). \\
\code{time\_marks} & $(r, t)$ pairs: fractional row and device time (Unix seconds) at the start and after every overflow. \\
\code{overflow\_marks} & (row, count) for each batch of UHD overflows. \\
\code{underflow\_marks} & (row, count) of transmit underflows, when the test transmitter is used. \\
\bottomrule
\end{longtable}

\begin{longtable}{@{}p{0.30\linewidth}p{0.65\linewidth}@{}}
\toprule
Attribute & Meaning \\
\midrule
\code{mode}, \code{observation\_mode} & \code{pulsar} \\
\code{dt\_s}, \code{nchan}, \code{channel\_width\_hz} & row length, channels, channel width \\
\code{sample\_rate\_hz}, \code{center\_freq\_hz}, \code{gain\_db}, \code{instrument} & the band as recorded \\
\code{t0\_unix}, \code{created\_utc} & host clock at start \\
\code{time\_marks\_exact} & 1 when the marks carry their exact fraction of a row (from 27 Sep 2026) \\
\code{time\_source}, \code{time\_pps\_verified}, \code{time\_device\_minus\_host\_s}, \code{time\_last\_pps} & how the absolute time was set \\
\code{sdr\_type}, \code{sdr\_serial} & the radio as UHD identifies it (\code{b200}) \\
\code{cal\_t\_sys\_k}, \code{cal\_effective\_area\_m2}, \code{cal\_t\_sys\_utc}, \code{cal\_beam\_utc} & the calibration in force at the start: $T_\mathrm{sys}$, $A_e$ and when each was measured (\code{pulsar\_fold.calibration\_in\_force}; from 30 Sep 2026) \\
\code{site\_lat\_deg}, \code{site\_lon\_deg}, \code{site\_height\_m} & 55.902426, $-4.307865$, 50\,m \\
\code{obs\_name}, \code{comment}, \code{object\_name}, \code{coord\_system}, \code{sun\_monitor}, \ldots & the schedule entry \\
\code{pulsar\_monitor}, \code{pulsar\_session} & a monitor run, and the window it belongs to (\S\ref{sec:monitor}); \code{pulsar\_session} is absent on a booking \\
\code{reference\_*} & the Thunderbolt's own account at the start: \code{reference\_state} (\code{absent} when no unit is read), lock, disciplining mode, holdover, alarms, oscillator and PPS offsets, DAC voltage, temperature, satellites \\
\code{pulsar\_name}, \code{pulsar\_period\_s}, \code{pulsar\_dm}, \ldots & the catalogue values used when the run was booked \\
\code{inject\_*} & present only for test-transmitter runs (\S\ref{sec:checks}) \\
\bottomrule
\end{longtable}

%-----------------------------------------------------------------------------
\section{Scheduling: the pulsar monitor}
\label{sec:monitor}

A pulsar run is either a booking (\code{coord\_system: pulsar}) or a run of the \textbf{pulsar monitor}, a standing order in \code{h1\_web\_scheduler.py} that records \psr{} once a day for \code{pulsar\_monitor\_hours} (16\,h). Its keys are \code{pulsar\_monitor} (off by default), \code{pulsar\_monitor\_window}, \code{pulsar\_monitor\_start} and \code{pulsar\_monitor\_hours}, on the Configuration tab.
\begin{description}
\item[The window.] \psr{} is circumpolar here (lowest 20.6$^\circ$, due north), but the measured horizon's northern obstruction, 50--55$^\circ$ from azimuth 340$^\circ$ to 360$^\circ$, hides it about 6\,h a day. That stretch comes 3.93\,min earlier each day, and moves an hour at each clock change, since the scheduler runs in local time.
  \begin{itemize}
  \item \code{follow} (the default): the window opens when the pulsar comes out from behind the active horizon profile, by the booking trim's full-beam test (\code{pulsar\_follow\_window}), and closes 16\,h later or when it goes behind again. Each day is one unbroken run in the $\sim$18\,h clear stretch: it opens at 19:42 in late September, 12:44 in late December and in the morning by March.
  \item \code{clock}: the window opens at \code{pulsar\_monitor\_start} every day and loses whatever the obstruction hides inside it. With no horizon profile measured there is no emergence to follow, and \code{follow} falls back to this.
  \end{itemize}
\item[Sessions.] A window is named by when it opens, to the minute (\code{pulsar\_session\_name}, e.g.\ \code{2026-09-28T1942}), so a window that crosses midnight is one session. The date alone is not unique: a following window moves 4\,min a day, so once a year two open on one date. Every run in the window carries the name as \code{pulsar\_session}.
\item[Precedence.] Bookings win. A monitor run ends 2\,min (\code{PULSAR\_MONITOR\_GUARD\_MINUTES}) before the next booking, and the monitor resumes afterwards for what is left of the window if at least 30\,min remain. Anything started by hand makes it give way through \code{sun\_monitor\_yield}, which holds both monitors off for 30\,min; \code{/api/stop} does the same. Over the Sun monitor it wins: no Sun run starts into the window, and one running when the window opens is stopped. A failed start is retried after 10\,min.
\item[Pieces.] A window interrupted by a booking or by hand comes back as several recordings under one session. With the PPS they share an absolute phase, so their folds add coherently into one night TOA (\S\ref{sec:toa}).
\item[Each run] is an ordinary pulsar-mode entry: homed first, the local horizon respected, stowed at the end.
\end{description}

When any pulsar run ends, booked or monitor, however it ended, \code{stop\_observation} hands the recording to \code{update\_pulsar\_timing}. That runs \code{pulsar\_toa.py} on it as a subprocess under \code{nice 19} and \code{ionice -c3}, since the next run may already be recording, and holds \code{observe\_plot\_lock}, because the plot writes the same files. It skips a recording under 10\,min (a failed start) or outside \code{data/observations/}, and logs every TOA line to \code{scheduler.log}. A 16\,h run takes about 25\,s.

%-----------------------------------------------------------------------------
\section{The phase model}
\label{sec:phase}

Every fold uses an \textbf{absolute} phase (\code{pulsar\_fold.absolute\_phase}), not one counted from the file's first sample. For each time $t$ (UTC at the site):
\begin{enumerate}
\item convert to TDB at the site's location (astropy);
\item add the light-travel time from the site to the solar-system barycentre, in the direction of the pulsar's position at that date, proper motion applied. This is the Roemer delay, up to 500\,s over the year;
\item subtract the dispersion delay $K\,\mathrm{DM}/f^2$ at the recording's centre frequency;
\item evaluate the spin ephemeris at that barycentric time, with $\delta = t_\mathrm{bary} - \mathrm{PEPOCH}$:
\begin{equation}
\phi = F_0\,\delta + \tfrac12 F_1\,\delta^2 + \tfrac16 F_2\,\delta^3 + \phi_\mathrm{off}.
\end{equation}
\end{enumerate}

The offset $\phi_\mathrm{off} = 0.3105$ was chosen once, from the first detection, so that its pulse folds to phase 0.5; with no offset that pulse fell at 0.1895. Every other run then lands at 0.5 \emph{by prediction}. Where a run actually lands is a combined check on the ephemeris, the barycentring and the clock. The offset must not be changed to centre any single night.

The phase is evaluated on 60\,s knots and interpolated linearly; the curvature between knots is below $10^{-7}$ of a period. Whole periods are removed at the run's start, so the numbers stay small: the $F_1$ term alone amounts to about 3300 periods by 2026. The fractional phase, which is what matters, is unaffected.

The older \code{phase\_track} integrated the topocentric period from the radial-velocity correction. It is kept only for the fixed-period test transmitter and for the \code{-f/-fd/-fdd} fit that prepfold needs.

\paragraph{Checks.}
\begin{itemize}
\item The rate $d\phi/dt$ agrees with the independent radial-velocity route to $5\times10^{-8}$. The remainder is TDB's rate against UTC, which the two routes treat differently.
\item Arrival times generated by PINT from the same par file, at our site, all fall at one fractional phase to within 1\,\textmu s over 5\,h.
\end{itemize}

%-----------------------------------------------------------------------------
\section{Folding}
\label{sec:fold}

\code{pulsar\_fold.analyse\_file} folds a recording in one pass with bounded memory. A 16\,h run of 57.5 million rows takes 24\,s and 260\,MB. Keeping whole-run arrays in memory had taken 10 minutes and 2.8\,GB.

\subsection{Per block (10 minutes of rows)}
\begin{enumerate}
\item Row times from the marks (mid-row), and phases interpolated from the absolute-phase knots.
\item \textbf{Baseline:} each channel is divided by its median over 1\,s blocks, less one. The result is a fractional excess over the system level, the same units at any bandwidth. The 1\,s block matches the 408\,MHz pipeline; a 1\% duty-cycle pulse does not move a median.
\item \textbf{Clipping:} samples beyond 6 robust $\sigma$ (1.4826\,MAD) are given zero weight, not set to zero. They are mostly positive-only broadband bursts.
\item \textbf{Weights:} every sample in a second gets the weight $1/\sigma^2$, where $\sigma^2$ is that second's variance over its unclipped samples. This is weighting by \emph{noise}, not by SNR. Weighting by measured SNR would favour stretches where the noise happened to line up with the pulse.
\item \textbf{Accumulate} $\sum w y$ and $\sum w$ into 60\,s sub-profiles of a 256-bin fold, and into a 1024-bin fold for timing when requested.
\end{enumerate}

\subsection{From the sub-profiles}
\begin{description}
\item[Profile.] The 256-bin fold summed over the run, and a 64-bin version for display.
\item[Matched SNR.] A boxcar of width $W_{50}$ run across the 256-bin profile; the peak over the median, in units of 1.4826\,MAD. This is the SNR quoted at the \emph{predicted} period, and the headline number.
\item[Period search.] The sub-profiles are shifted by whole fine bins, as a period offset $d$ would move them, and re-summed for $\pm150$\,ppm. The search is a check, not the quoted number: pure noise reaches SNR $\sim$3.5 somewhere in the range. For the first detection it peaked at $+0.0$\,ppm.
\item[Sub-integrations.] 10-minute slices at the predicted period, for the time--phase image.
\end{description}

The quoted SNR moves by a few tenths depending on where the pulse falls against the bin edges. The second decimal place of any single SNR is not meaningful.

\subsection{Expected sensitivity}
\label{sec:expect}
In the fold's fractional units, noise per 1\,ms row is $1/\sqrt{B\,\Delta t}$: 0.0112 at 8\,MHz and 0.0060 at 28\,MHz effective. Measured values are 1.05 and 1.08 times these. For a Gaussian pulse, the matched-filter SNR over an integration $T$ is
\[
\mathrm{SNR} \propto \frac{S\,G}{T_\mathrm{sys}}\sqrt{B\,T}\ \times\ \eta_\mathrm{profile}\,\eta_\mathrm{pol},
\]
with $G = A_e/2k = 2.55$\,mK\,Jy$^{-1}$ for one polarisation ($A_e = \lambda^2/\Omega_\mathrm{main} = 7.05$\,m$^2$) and $T_\mathrm{sys} \approx 187$\,K. Two factors reduce it:
\begin{itemize}
\item $\eta_\mathrm{profile} = 0.59$: the real profile holds less of its energy in the core than a single 6.6\,ms Gaussian.
\item $\eta_\mathrm{pol} = 0.63$--$1.4$: the fraction of the circularly polarised core the single feed takes, which depends on the feed's unknown handedness.
\end{itemize}
The result is about 3.5$\sigma$ for 4\,h at 8\,MHz, and about 13$\sigma$ for 16\,h at 32\,MHz, 8--19$\sigma$ with the polarisation factor. Scintillation at 1.4\,GHz is coarse (scintles of order MHz, lasting minutes to tens of minutes), so a 28\,MHz band averages only a handful. A night can be twice or half the average.

\subsection{Results so far}
\begin{longtable}{@{}p{0.34\linewidth}p{0.61\linewidth}@{}}
\toprule
Run & Result \\
\midrule
26 Sep 01:00, 4\,h, 8\,MHz & matched SNR 2.0: not detected; 0 overflows; through transit \\
26 Sep 08:00, 4\,h, 8\,MHz & about 3: not detected \\
26--27 Sep, 16\,h, 32\,MHz & matched SNR 9.4--9.5, template SNR 14.9; period search +0.0\,ppm; PRESTO $9.9\sigma$; halves 7.8 and 7.2 at the same phase \\
27--28 Sep, 16\,h, 32\,MHz, host clock & template SNR 15.7 \\
28--29 Sep, 16\,h, 32\,MHz, PPS & first monitor session, \code{2026-09-28T1942}, one piece; template SNR 9.5 \\
29--30 Sep, 15.5\,h, 32\,MHz, PPS & session \code{2026-09-29T1938} in two pieces, 19:39--06:25 and 06:57--11:39; night template SNR 14.0 \\
\bottomrule
\end{longtable}

\paragraph{Flux.} The flux density follows from the period-averaged pulse area over $\pm30$\,ms (so the outriders are included), times $T_\mathrm{sys}/G$ (\code{pulsar\_fold.mean\_flux\_mjy}). The scale is each recording's own: from 30 September the receiver writes the calibration in force at the start into the file (\code{cal\_*}, \S\ref{sec:file}), so a later re-calibration never rescales an old night. An earlier file falls back to the calibration in force when it is plotted, and the plot says \code{current} rather than \code{recorded}. Results:
\begin{itemize}
\item first detection: $188\pm17$\,mJy;
\item night of 29--30 September: $184\pm21$\,mJy;
\item all runs stacked, 71\,h: $168\pm8$\,mJy;
\item ATNF: $203\pm57$\,mJy.
\end{itemize}
The errors are statistical. Systematic uncertainty is about $\pm40\%$: $T_\mathrm{sys}$ is the H\,\textsc{i} band's, not measured separately for the 32\,MHz pulsar band; the single feed's share of a pulse with $V/I = -0.34$ is unknown; and pointing.

\paragraph{Scintillation.} On the first night the pulsar brightened as the run went on: $135\pm24$ then $242\pm25$\,mJy by halves, and 91--368\,mJy in 2\,h pieces. The altitude count drift found at the next homing (0.5\,deg beyond the usual reading) would have lowered the late signal, not raised it.

%-----------------------------------------------------------------------------
\section{Times of arrival}
\label{sec:toa}

\subsection{Template}
The template is the EPN 1410\,MHz Stokes-$I$ profile (hx97b, Effelsberg, 4096 bins), stored in \code{pulsar\_templates/}. Because it is noise-free and independent of our data, it cannot pull a TOA towards our own noise. It is smoothed exactly as our system smooths:
\begin{itemize}
\item one row (1\,ms);
\item the dispersion sweep across the recording's band;
\item one fold bin (1/1024 of a period).
\end{itemize}
Each is applied as a boxcar through its transfer function, a sinc. The main peak is placed at phase 0.5 and the template normalised to a peak of one.

Our single feed sees $I$ plus a mixture of $Q$, $U$ and $V$, so our profile shape differs slightly from $I$. A constant difference becomes a phase offset, which the timing fit absorbs. One that changes with parallactic angle through a track would show as TOAs that depend on hour angle, and has to be tested for.

\subsection{Fit}
The run's 1024-bin profile $p$ is fitted as
\begin{equation}
p(\phi) = a + b\,T(\phi - \tau)
\end{equation}
in the Fourier domain (Taylor 1992), in \code{pulsar\_toa.fit\_shift}. The mean harmonic is excluded, so $a$ drops out.
\begin{enumerate}
\item A 16-times oversampled cross-correlation finds the peak approximately.
\item A bounded one-dimensional maximisation of $\mathrm{Re}\sum_k P_k T_k^* e^{2\pi i k\tau}$ refines it, so $\tau$ is continuous rather than stepped by a bin.
\item $b$ is fitted by least squares at that $\tau$.
\end{enumerate}
The error comes from the Fisher information:
\begin{equation}
\sigma_\tau = \frac{\sigma_\mathrm{bin}}{|b|\,\sqrt{\sum_j T'(\phi_j-\tau)^2}},
\end{equation}
with $\sigma_\mathrm{bin}$ the residual scatter per bin. The template SNR is $b\sqrt{\sum T^2}/\sigma_\mathrm{bin}$.

For a resolved pulse the bin width drops out: splitting finer raises the per-bin noise by exactly as much as it sharpens the shape. For a Gaussian the precision is $\sigma_t \approx 0.6\,\mathrm{FWHM}/\mathrm{SNR}$, about 0.4\,ms for a 15$\sigma$ night.

In 300 Monte Carlo trials the fitted shift was unbiased and its scatter matched the quoted error ($1.00\pm0.15$). On the first detection, $\sigma_\mathrm{bin}$ was 0.96 times the radiometer prediction.

\subsection{Arrival time}
The TOA is the time at which the folding phase model reads $n + 0.5 + \tau$, taking the pulse nearest the middle of the span. It is found by interpolating the phase knots (\code{arrival\_time}). This is a \textbf{topocentric UTC} time at the site and the \textbf{observing frequency}; barycentring and dispersion are left to the timing program, which applies them from the par file. Its error is $\sigma_\tau P_\mathrm{topo}$. The MJD is written to 13 decimal places from astropy's two-part time, never through a single float64, whose MJD resolution is about 0.6\,\textmu s.

A TOA is written only when the template SNR is at least 6 (\code{MIN\_TOA\_SNR}); below that a fit can lock onto a noise feature.
\begin{itemize}
\item \textbf{Night TOA:} one per night. This is the timing datum, and the only kind fitted.
  \begin{itemize}
  \item A booked run gives \code{<recording>}, from the profile of the whole run folded at once.
  \item A pulsar-monitor session gives \code{night\_<session>} (\code{pulsar\_toa.session\_toa}): the stored profiles of all the session's pieces added at the absolute phase and fitted as one. It sits at the pieces' duration-weighted middle, its \code{-tobs} is the time recorded rather than the span, and it carries \code{-pps 1} only if every piece was timed by the PPS. It is recomputed as each piece arrives, so each piece's profile is stored first. Each piece's own TOA, \code{<recording>\_p}, goes to the segments file as a check.
  \end{itemize}
\item \textbf{Segments} (\code{<recording>\_sN}): equal parts of a run, each fitted separately. They are still written for every run, monitor pieces included. The scheduler's end-of-run update cuts a run into $\mathrm{round}(h/4)$ parts for a run of $h$ hours (\code{PULSAR\_SEGMENT\_HOURS}), so 16\,h gives four and the 10.8\,h first piece of 29--30 September gave three; Plot Result cuts into four. Rounding rather than truncating matters: a 16\,h window records 15.98\,h, and truncation once gave three 5.3\,h segments that overwrote the 4\,h lines of the same names. A run's segments share its clock offset, so their differences are exact even without the PPS. They are a within-night check, never fitted. The first night's four segments (SNR 6.1, 6.9, 6.4, 10.9) scatter at $\chi^2 = 10.9$ for 3 degrees of freedom. Three of the four sit near the SNR cut, where fits are least reliable.
\end{itemize}

\subsection{The .tim and par files}
Each TOA is one line in tempo2 format. Night TOAs go in \code{pulsar\_timing/b0329.tim}, and segments and monitor pieces in \code{pulsar\_timing/b0329\_segments.tim}; \code{write\_tim} routes each line by its name. The two are kept apart because a segment or a piece is the same data as its night again: a timing program given both in one file cannot know that, and would count every night twice. \code{b0329.tim} is the file to fit, with any timing program; the segments file is a check on it and is never fitted with it. The line format is the same in both:
\begin{verbatim}
<name> <freq MHz> <MJD UTC> <error us> ar -pps 0|1 -snr S -tobs s -bw MHz
       -be b200 -fe sawbird -tmpl epn_hx97b_1410
\end{verbatim}
Lines are keyed by name, so re-analysing a recording replaces its lines rather than duplicating them, and the file is kept in time order. \code{-pps 0} marks a host-clock TOA (the first two nights), \code{-pps 1} one timed by the PPS (every night from 28 September); the par file adds \code{EQUAD -pps 0 2000} (\textmu s) in quadrature for those. \code{B0329+54.par} is written from the same parameters the fold uses (\code{UNITS TDB}, \code{EPHEM DE421}, \code{CLK TT(TAI)}). A JSON record next to each line keeps the full fit.

%-----------------------------------------------------------------------------
\section{Timing with PINT}
PINT runs in its own conda environment, \code{radioconda/envs/pint}, never in the observing environment. \code{pint\_tools.py} is driven by subprocess and exchanges JSON.
\begin{itemize}
\item \textbf{Site.} Our surveyed position is registered as a custom observatory, \code{acre\_srt}, alias \code{ar}. PINT's own \code{acre} lies 49\,m away. No GPS or site clock corrections are applied, because the times are UTC from NTP or the PPS.
\item \textbf{Fit.} Night TOAs only, from \code{b0329.tim}. The segments get residuals under the fitted model, for the plot. PEPOCH is moved to the mean epoch of the data, so $P$ and $\dot P$ are quoted there rather than extrapolated 40 years. $F_0$ is freed once three night TOAs span at least a day. $F_1$ is freed once four span three weeks: the spin-down signature grows as $\tfrac12 F_1 t^2/F_0$, about 5\,ms in three weeks and 10\,ms in a month, against TOA errors of 0.4--2\,ms. Before that, the parameters are held at the catalogue values and only residuals are shown.
\item \textbf{Checks.} PINT reads our \code{.tim} back with residuals below 0.1\,\textmu s. Our absolute phase and PINT's arrivals agree to 1\,\textmu s over 5\,h.
\item \textbf{So far} (fit of 30 September). Four night TOAs over 2.96 days, template SNR 14.9, 15.7, 9.5 and 14.0, the last two on the PPS. $F_0$ free: $\chi^2 = 1.00$ on 2 degrees of freedom, rms 0.73\,ms, $P = 0.7145223292(39)$\,s at the data epoch. The catalogue ephemeris gives 0.7145223255\,s there, $0.9\sigma$ away. $F_1$ frees itself when four nights span three weeks; recovering $\dot P$ is the outstanding check. Issue \#48 (absolute TOAs on the PPS) was closed on 30 September.
\end{itemize}

%-----------------------------------------------------------------------------
\section{PRESTO as an independent check}
PRESTO is installed in its own environment, and its \code{prepfold} is rebuilt with \code{tools/presto\_acre\_road.patch} so that it knows our site: telescope id 82, observatory code AR.

\code{sigproc\_export.py} writes a SIGPROC \code{.fil} file. Channels are written highest first, as 32-bit floats. Every overflow gap is filled with rows at the channel median, because a \code{.fil} has one time grid: sample $i$ is at $t_\mathrm{start} + i\,t_\mathrm{samp}$. The export therefore rounds each post-overflow stretch to whole rows.

\code{presto\_fold} runs
\begin{verbatim}
prepfold -topo -f f -fd fd -fdd fdd -phs phi0 -dm 26.7641 -n 64
         -nsub 1 -nosearch -scaleparts
\end{verbatim}
Here $(f, \dot f, \ddot f)$ come from a cubic fit to our absolute phase over the run, and $-\mathrm{phs}$ is our absolute phase at the first sample, so PRESTO's pulse also sits at 0.5. The DM is never searched: 32\,MHz at 1.4\,GHz does not constrain it. PRESTO's significance is a reduced $\chi^2$ of the profile against a flat line, a weaker test than a matched filter. Each tool's number is quoted as its own; they are not interchangeable.

%-----------------------------------------------------------------------------
\section{The pulsar plot}
Plot Result on a pulsar recording runs one fold with the timing bins and draws four panels.
\begin{description}
\item[Top left] This run's profile at the predicted period, in blue, with the best-searched period in orange. Two periods, pulse at 0.5 and 1.5; baseline at the off-pulse median. Drawn in antenna temperature, mK, as the fraction of total power times $T_\mathrm{sys}$ (\code{radiometer\_scale}), with \% of total power as a secondary axis on the right; in \% alone when there is no calibration. The title carries the matched SNR and the period-averaged flux density, with the $T_\mathrm{sys}$ and $A_e$ used and whether they are the recording's own (\code{recorded}) or the calibration in force now (\code{current}).
\item[Bottom left] 10-minute sub-integrations on exactly the same two-period axis: the pixel edges fall on the profile's bin edges, and the axes are shared.
\item[Top right] All runs stacked at the absolute phase from the stored profiles, phase 0.25--0.75, with the fitted template overlaid. Each night is converted to kelvin on its own $T_\mathrm{sys}$ before the nights are added, with the same noise weights (\code{accumulated\_profile(calibrated=True)}), so a stack of nights under different calibrations is still on one scale. It is drawn in mK with its flux density in a box, and the template at the stack's weighted-mean $T_\mathrm{sys}$. If any night has no calibration, own or current, it falls back to units of the fitted template's peak.
\item[Bottom right] PINT residuals: night TOAs in red (filled circles for PPS time, open squares for host clock), segments and pieces in grey, statistical errors only, with $P$ and $\dot P$ fitted or held at the data epoch. A host-clock night also gets a faint outer bar, statistical and the 2\,ms EQUAD together, which is what the fit weights it by.
\end{description}

For a file in \code{data/observations/}, plotting also writes the run's TOAs and its stored profile, as the end-of-run update does (\S\ref{sec:monitor}). A test's temporary file never touches either. Plot and PRESTO requests are serialised by \code{observe\_plot\_lock}: pyplot is not thread-safe, and two plots share their outputs.

%-----------------------------------------------------------------------------
\section{Where the data are held}
\begin{longtable}{@{}p{0.38\linewidth}p{0.12\linewidth}p{0.44\linewidth}@{}}
\toprule
Location & In git & Contents \\
\midrule
\code{data/observations/*\_pulsar.h5} & no & raw recordings (\S\ref{sec:file}), 14\,MB per hour; re-observable \\
\code{pulsar\_timing/b0329.tim} & yes & night TOAs, one per booked run or monitor session: the file to fit \\
\code{pulsar\_timing/b0329\_segments.tim} & yes & segment TOAs (about 4\,h) and monitor pieces, a check; never fitted with the nights \\
\code{pulsar\_timing/B0329+54.par} & yes & the ephemeris the fold and PINT start from \\
\code{pulsar\_timing/<name>.json} & yes & full record behind each TOA line \\
\code{pulsar\_timing/profiles/<rec>.npz} & yes & each run's whole-run fold: 1024-bin $\sum wy$ and $\sum w$ at the absolute phase, with \code{phase\_offset}, \code{row\_centre}, \code{nbins}, frequency, bandwidth, time span, clock source, \code{pps}, \code{session}, and the calibration: \code{cal\_t\_sys\_k}, \code{cal\_effective\_area\_m2}, and \code{cal\_recorded} (1 if taken from the recording, 0 if from the calibration in force when the profile was stored) \\
\code{pulsar\_templates/} & yes & EPN template and its source file (CC BY 4.0, attribution in README) \\
\code{data/presto/} & no & \code{.fil} exports and prepfold output \\
\bottomrule
\end{longtable}

The TOAs and stored profiles are kept in git because each represents a night of telescope time, and a TOA history cannot be rebuilt from any single recording. The stacked profile is built from the stored profiles, not from the recordings on disk, so a night still counts after its recording has gone. A profile made under a different \code{phase\_offset} is rotated to the current one by a Fourier shift. The end-of-run update and Plot Result write these files; committing them is done by hand.

%-----------------------------------------------------------------------------
\section{How each step was established}
\label{sec:checks}
\begin{description}
\item[Artificial pulsar.] The B200's transmitter sends gated noise through a 30\,dB pad and the vertex dipole (\code{H1\_PULSAR\_INJECT}). The replica mode shapes it as a 6.6\,ms Gaussian at the topocentric phase given by our own model.
  \begin{itemize}
  \item All 439 single pulses arrived at a fixed phase (MAD 0.66\,ms).
  \item Pulse area was within 6\% of the level check, and scaled as amplitude$^2$ between amplitudes 0.1 and 0.2.
  \item Fold noise was 0.96--1.12 times the radiometer prediction.
  \end{itemize}
  Below amplitude 0.05 the injection is unusable: whenever the transmit samples are non-zero, the received power dips by a roughly fixed amount. A transmit underflow slips the whole pulse train by 29--61\,ms, so the injection chain carries about a second of buffering.
\item[Synthetic data.] Buried pulses are recovered at their phase. Two runs a day apart fold to the same phase. A run deliberately offset from the ephemeris returns the offset through its TOA. Error bars are checked by Monte Carlo. Overflow marks are exact to the sample.
\item[PINT.] The phase agreement and \code{.tim} round trip described above.
\item[The sky.] The detection: period search at $+0.0$\,ppm, both halves at the same phase, PRESTO's independent fold, and a flux consistent with the catalogue.
\end{description}

%-----------------------------------------------------------------------------
\section{Open questions and next steps}
\begin{itemize}
\item \textbf{PPS.} Wired and verified: every run from 28 September has \code{time\_source} \code{pps} with \code{time\_pps\_verified} 1, and its night TOA carries \code{-pps 1}. Only the first two nights carry the host clock's 1--4\,ms, and the fit weights them by the 2\,ms EQUAD.
\item \textbf{Segment scatter} on the first night ($\chi^2 = 10.9/3$). Two candidates: that file's $\pm0.5$\,ms overflow rounding (it predates exact marks), and the feed's view of the polarisation turning with parallactic angle; the $+2$\,ms segment contains transit. Weak fits near the SNR cut are a third. The test is TOAs split by hour angle across several nights.
\item \textbf{Zenith pointing} (issue \#49). The controller clamps the model's azimuth correction at 10$^\circ$, which truncates the collimation term above altitude $\approx82^\circ$. On the transit of \psr{} the beam falls up to 1.7$^\circ$ off-target, costing about 2\% of a 16\,h run's signal. The fix belongs in the firmware.
\item \textbf{Count drift.} About one pulse (0.5$^\circ$) per axis in 16--18\,h of tracking. Every monitor run homes first.
\item \textbf{$\dot P$.} Four nights now span three days; $F_1$ is freed when four span three weeks, and recovering the catalogue $\dot P$ is the next check on the whole chain.
\item \textbf{Interference.} A broadband band at about 640\,min into the first detection run lights every phase at once. Lines at multiples of 12.5\,Hz and red noise at 0.05--0.3\,Hz are present on the sky and at the stow, but cost the fold only about 10\%.
\item \textbf{Sensitivity.} A system temperature of 100\,K would give about 30$\sigma$ in 16\,h, or 5$\sigma$ in about 30 minutes. A second polarisation would add $\sqrt2$ and remove the feed-handedness uncertainty, but the B200 has one receive channel: it needs a B210 or a second B200 on the same reference.
\end{itemize}

%-----------------------------------------------------------------------------
\section*{References}
\begin{itemize}
\item Hobbs, G. et al.\ 2004, MNRAS 353, 1311: timing noise and the ATNF parameters of \psr{}.
\item Manchester, R.\,N., Hobbs, G.\,B., Teoh, A., Hobbs, M.\ 2005, AJ 129, 1993: the ATNF pulsar catalogue.
\item Taylor, J.\,H.\ 1992, Phil.\ Trans.\ R.\ Soc.\ A 341, 117: pulsar timing and template matching in the Fourier domain.
\item von Hoensbroech, A. \& Xilouris, K.\,M.\ 1997, A\&AS 126, 121: Effelsberg multifrequency pulsar polarimetry (EPN dataset hx97b).
\item Luo, J. et al.\ 2021, ApJ 911, 45: PINT.
\item Ransom, S.\,M.\ 2001, PhD thesis, Harvard: PRESTO.
\end{itemize}

\end{document}
